\section{Conclusion}
We explored representation engineering (RepE), an approach to top-down transparency for AI systems. Inspired by the Hopfieldian view in cognitive neuroscience, RepE places representations and the transformations between them at the center of analysis. As neural networks exhibit more coherent internal structures, we believe analyzing them at the representation level can yield new insights, aiding in effective monitoring and control. Taking early steps in this direction, we proposed new RepE methods, which obtained state-of-the-art on TruthfulQA, and we demonstrated how RepE and can provide traction on a wide variety of safety-relevant problems. While we mainly analyzed subspaces of representations, future work could investigate trajectories, manifolds, and state-spaces of representations. We hope this initial step in exploring the potential of RepE helps to foster new insights into understanding and controlling AI systems, ultimately ensuring that future AI systems are trustworthy and safe.
